% Part: turing-machines
% Chapter: machines-computations
% Section: combining-machines

\documentclass[../../../include/open-logic-section]{subfiles}

\begin{document}

\olfileid[hi]{tur}{mac}{cmb}
\olsection{ट्यूरिंग मशीनों का संयोजन}

\begin{explain}
अब तक देखे गए ट्यूरिंग मशीनों के उदाहरण स्वभावतः अपेक्षाकृत सरल रहे
हैं। किंतु वास्तव में आधुनिक प्रोग्रामन की किसी भी भाषा से हल की जा सकने
वाली प्रत्येक समस्या ट्यूरिंग मशीनों से भी हल की जा सकती है। अधिक जटिल
ट्यूरिंग मशीनें बनाने के लिए यह स्थापित करना महत्त्वपूर्ण है कि हम उनका
संयोजन कर सकते हैं, ताकि प्रक्रिया को सरल भागों में बाँटकर अधिक जटिल
समस्याएँ हल करने वाली मशीनें बना सकें। यदि किसी जटिल समस्या को उसके घटक
भागों में बाँटने की स्वाभाविक रीति मिल जाए, तो कई सरल ट्यूरिंग मशीनें
बनाकर और उन्हें समस्या हल करने वाली एक मशीन में मिलाकर समस्या को कई
चरणों में हल किया जा सकता है। अगले अनुभाग में विराम समस्या पर काम करते
समय यह बात विशेष रूप से महत्त्वपूर्ण होगी।

ट्यूरिंग मशीनों $M = \tuple{Q, \Sigma, q_0, \delta}$ और
~$M' = \tuple{Q', \Sigma', q_0', \delta'}$ का संयोजन कैसे करें? अब हम
$M$ के रुकने के बाद पट्टी के विन्यास को मशीन~$M'$ की किसी क्रिया के आगत
विन्यास के रूप में लेते हैं। यह करने वाली एकल ट्यूरिंग मशीन
$M \frown M'$ पाने के लिए निम्न कार्य करें:
\begin{enumerate}
    \item $Q'$ की सभी अवस्थाओं को, जो~$M'$ की हैं, पुनः संख्या (या पुनः नाम) दें,
    ताकि $M$ और~$M'$ की कोई अवस्था समान न हो
    ($Q \cap Q' = \emptyset$)।
    \item $M \frown M'$ की अवस्थाएँ $Q \cup Q'$ हैं।
    \item पट्टी की वर्णमाला $\Sigma \cup \Sigma'$ है।
    \item आरंभिक अवस्था~$q_0$ है।
    \item संक्रमण फलन निम्न फलन $\delta''$ है:
    \[\delta''(q,\sigma) =
    \begin{cases}
      \delta(q,\sigma) & \text{यदि $q \in Q$}\\
      \delta'(q,\sigma) & \text{यदि $q \in Q'$}\\
      \tuple{q_0', \sigma, \TMstay} & \text{यदि $q \in Q$ और
      $\delta(q,\sigma)$ अपरिभाषित है}
    \end{cases}\]
\end{enumerate}
परिणामी मशीन~$M$ के निर्देशों का उपयोग तब करती है जब वह अवस्था $q \in Q$
में हो, और~$M'$ के निर्देशों का उपयोग तब करती है जब वह अवस्था~$q \in
Q'$ में हो। जब वह अवस्था
$q \in Q$ में है और ऐसा प्रतीक~$\sigma$ पढ़ रही है जिसके लिए $M$ में
कोई संक्रमण नहीं है (अर्थात् $M$ रुक जाती), तब वह~$M'$ की आरंभिक अवस्था
में प्रवेश करती है (और पट्टी की सामग्री तथा शीर्ष की स्थिति को यथावत्
छोड़ती है)।

ध्यान दें कि यदि मशीन~$M$ अनुशासित नहीं है, तो $M$ के रुकते समय हमें
नहीं पता कि पट्टी का शीर्ष कहाँ है; अतः आवश्यक नहीं कि~$M$ के अंतिम
विन्यास में शीर्ष खाना~$1$ पढ़ रहा हो। मशीनों का संयोजन करते समय इसे
ध्यान में रखना महत्त्वपूर्ण है।
\end{explain}

\begin{ex}
\emph{मशीनों का संयोजन:} हम ऐसी मशीन की अभिकल्पना करेंगे जो~$\TMstroke$
के दो खंडों के आगत पर आरंभ हो, जिनकी लंबाइयाँ क्रमशः $n$ और~$m$ हों, और पट्टी
पर $2(m+n)$ $\TMstroke$ के एकल खंड के साथ रुके। इस मशीन को बनाने के लिए
हम दो परिचित मशीनों का संयोजन कर सकते हैं: योग मशीन और द्विगुणक। योग
मशीन का अवस्था आरेख बनाकर आरंभ करें।
\[
\begin{tikzpicture}[->,>=stealth',shorten >=1pt,auto,node distance=2.8cm,
                    semithick]
  \tikzstyle{every state}=[fill=none,draw=black,text=black]

  \node[initial,state] (A)              {$q_0$};
  \node[state]         (B) [right of=A] {$q_1$};
  \node[state]         (C) [right of=B] {$q_2$};

  \path (A) edge node {\TMtrans{\TMblank}{\TMstroke}{\TMstay}} (B)
            edge [loop above] node {\TMtrans{\TMstroke}{\TMstroke}{\TMright}} (B)
        (B) edge [loop above] node {\TMtrans{\TMstroke}{\TMstroke}{\TMright}} (B)
            edge node {\TMtrans{\TMblank}{\TMblank}{\TMleft}} (C)
        (C) edge [loop above] node {\TMtrans{\TMstroke}{\TMblank}{\TMstay}} (C);
\end{tikzpicture}
\]
अवस्था~$q_2$ में रुकने के स्थान पर निर्गत को दुगुना करने के लिए क्रिया
जारी रखना चाहते हैं। याद करें कि द्विगुणक मशीन आगत का पहला आघात मिटाकर
अलग निर्गत में दो आघात लिखती है। यह सुनिश्चित करने वाला निर्देश जोड़ें
कि पट्टी का शीर्ष योग मशीन के निर्गत का पहला आघात पढ़ रहा हो।
\[
\begin{tikzpicture}[->,>=stealth',shorten >=1pt,auto,node distance=2.8cm,
                    semithick]
  \tikzstyle{every state}=[fill=none,draw=black,text=black]

  \node[initial,state] (A)              {$q_0$};
  \node[state]         (B) [right of=A] {$q_1$};
  \node[state]         (C) [right of=B] {$q_2$};
  \node[state]         (D) [below left of=C] {$q_3$};
  \node[state]         (E) [below left of=D] {$q_4$};

  \path (A) edge node {\TMtrans{\TMblank}{\TMstroke}{\TMstay}} (B)
            edge [loop above] node {\TMtrans{\TMstroke}{\TMstroke}{\TMright}} (B)
        (B) edge [loop above] node {\TMtrans{\TMstroke}{\TMstroke}{\TMright}} (B)
            edge node {\TMtrans{\TMblank}{\TMblank}{\TMleft}} (C)
        (C) edge node {\TMtrans{\TMstroke}{\TMblank}{\TMleft}} (D)
        (D) edge [loop left] node {\TMtrans{\TMstroke}{\TMstroke}{\TMleft}} (D)
            edge node[left, xshift=-2mm] {\TMtrans{\TMendtape}{\TMendtape}{\TMright}} (E);
\end{tikzpicture}
\]
अब आगत को दुगुना करना सरल है---हमें केवल द्विगुणक मशीन को अवस्था~$q_4$
पर जोड़ना है। इसके लिए द्विगुणक मशीन की अवस्थाओं का पुनर्नामकरण करना
होगा, ताकि वे~$q_4$ से आरंभ हों, न कि~$q_0$ से---इस प्रकार दो
आरंभिक अवस्थाएँ नहीं बनेंगी। अंतिम आरेख \olref{fig:combined} जैसा दिखना
चाहिए।
\begin{figure}
\[
\begin{tikzpicture}[->,>=stealth',shorten >=1pt,auto,node distance=2.8cm,
                    semithick]
  \tikzstyle{every state}=[fill=none,draw=black,text=black]
  \node[initial,state] (A)              {$q_0$};
  \node[state]         (B) [right of=A] {$q_1$};
  \node[state]         (C) [right of=B] {$q_2$};
  \node[state]         (D) [below left of=C] {$q_3$};
  \node[state]         (E) [below left of=D] {$q_4$};
  \node[state]         (2) [right of=E] {$q_5$};
  \node[state]         (3) [right of=2] {$q_6$};
  \node[state]         (4) [below of=3] {$q_7$};
  \node[state]         (5) [left of=4]  {$q_8$};
  \node[state]         (6) [left of=5]  {$q_9$};

  \path (A) edge node {\TMtrans{\TMblank}{\TMstroke}{\TMstay}} (B)
            edge [loop above] node {\TMtrans{\TMstroke}{\TMstroke}{\TMright}} (B)
        (B) edge [loop above] node {\TMtrans{\TMstroke}{\TMstroke}{\TMright}} (B)
            edge node {\TMtrans{\TMblank}{\TMblank}{\TMleft}} (C)
        (C) edge node {\TMtrans{\TMstroke}{\TMblank}{\TMleft}} (D)
        (D) edge [loop left] node {\TMtrans{\TMstroke}{\TMstroke}{\TMleft}} (D)
            edge node[left, xshift=-2mm] {\TMtrans{\TMendtape}{\TMendtape}{\TMright}} (E)
    (E) edge node {\TMtrans{\TMstroke}{\TMblank}{\TMright}} (2)
    (2) edge [loop above] node {\TMtrans{\TMstroke}{\TMstroke}{\TMright}} (2)
      edge node {\TMtrans{\TMblank}{\TMblank}{\TMright}} (3)
    (3) edge [loop above] node {\TMtrans{\TMstroke}{\TMstroke}{\TMright}} (3)
        edge node {\TMtrans{\TMblank}{\TMstroke}{\TMright}} (4)
    (4) edge [loop below] node {\TMtrans{\TMblank}{\TMstroke}{\TMleft}} (4)
        edge node {\TMtrans{\TMstroke}{\TMstroke}{\TMleft}} (5)
    (5) edge [loop below]  node {\TMtrans{\TMstroke}{\TMstroke}{\TMleft}} (5)
        edge              node {\TMtrans{\TMblank}{\TMblank}{\TMleft}} (6)
    (6) edge [loop below] node {\TMtrans{\TMstroke}{\TMstroke}{\TMleft}} (6)
        edge              node {\TMtrans{\TMblank}{\TMblank}{\TMright}} (E);
\end{tikzpicture}
\]
\caption{योग और द्विगुणक मशीनों का संयोजन}
\ollabel{fig:combined}
\end{figure}
\end{ex}

\begin{prop}
    यदि $M$ और $M'$ अनुशासित हैं और क्रमशः फलनों $f\colon
    \Nat^k \to \Nat$ और $f'\colon \Nat \to \Nat$ को संगणित करती हैं,
    तो $M \frown M'$ अनुशासित है और~$\comp{f}{f'}$ को संगणित करती है।
\end{prop}

\begin{proof}
    चूँकि $M$ अनुशासित है, इसलिए जब वह निर्गत
    ~$f(n_1,\dots,n_k) = m$ के साथ रुकती है, तब शीर्ष खाना~$1$ पढ़ रहा
    होता है। अब यदि हम~$M'$ की आरंभिक अवस्था में प्रवेश करें, तो मशीन
    निर्गत $f'(m)$ के साथ और फिर से खाना~$1$ पढ़ते हुए रुकेगी।
    \olref[dis]{defn:disciplined} की अन्य शर्तें भी पूरी होती हैं।
\end{proof}

\begin{prob}
    $f(x) = x+2$ संगणित करने के लिए मशीन~$M$ को
    \cref{tur:mac:dis:prob:disc-succ} से लें और $M \frown M$ बनाएँ; इस
    प्रकार प्राप्त अनुशासित ट्यूरिंग मशीन दें।
\end{prob}
\end{document}
